% Diagrama do protocolo: as duas fases compartilham a mesma cauda determinística,
% que é a afirmação central do capítulo. Estágio probabilístico com borda
% tracejada, determinístico com borda sólida: a distinção é o conteúdo da figura,
% não decoração, e sobrevive à impressão em tons de cinza.
%
% Escrito com coordenadas explícitas e sem `\resizebox`, como as figuras BPMN de
% figuras/tikz/. A primeira versão usava `positioning` (`right=of`) e envolvia a
% imagem em `\resizebox`; com o preâmbulo dessincronizado, a chave desconhecida
% quebrava o `tikzpicture` dentro da caixa e derrubava o TikZ de TODAS as figuras
% seguintes. Aqui não há dependência além de `arrows.meta`, já usada pelas
% figuras BPMN, e a largura cabe no texto por construção.
\begin{figure}[htb]
	\Caption{\label{fig:pipeline} Visão geral: construção do corpus e avaliação}
	\centering
	\begin{tikzpicture}[
		x=1cm, y=1cm, font=\footnotesize,
		caixa/.style={draw, rounded corners=2pt, align=center, text width=2.0cm, minimum height=1.1cm, inner sep=2pt, fill=white},
		prob/.style={caixa, dashed, thick, draw=corDSL},
		det/.style={caixa, thick, draw=corXML},
		art/.style={align=center, font=\scriptsize\itshape, text=black!70},
		papel/.style={align=center, font=\scriptsize, text=black!60},
		fase/.style={anchor=west, font=\footnotesize\bfseries, text=black!70},
		seta/.style={-{Latex[length=2mm]}, thick, draw=black!55}]
		% Faixa das transformações determinísticas, comum às duas fases.
		\fill[black!5, rounded corners=3pt] (7.2,-6.1) rectangle (15.75,1.45);
		\node[papel] at (11.45,1.2) {transformações determinísticas};
		% Construção do corpus.
		\node[fase] at (-0.1,1.75) {Construção do corpus};
		\node[art] (txt) at (0.4,0) {texto\\bruto};
		\node[prob] (pre) at (2.3,0) {Padronização\\do texto\\\scriptsize (LLM)};
		\node[prob] (ext) at (5.5,0) {Extração\\para JSON\\\scriptsize (LLM)};
		\node[det] (conv) at (8.6,0) {Conversão\\JSON para DSL};
		\node[det] (trans) at (11.45,0) {Transpilação\\DSL para XML};
		\node[det] (val) at (14.4,0) {Validação\\\scriptsize XSD};
		\draw[seta] (txt) -- (pre);
		\draw[seta] (pre) -- (ext);
		\draw[seta] (ext) -- (conv);
		\draw[seta] (conv) -- (trans);
		\draw[seta] (trans) -- (val);
		\node[art] at (3.9,1.08) {texto\\padronizado};
		\node[art] at (7.05,0.8) {JSON};
		\node[art] at (10.05,0.8) {DSL};
		\node[art] at (12.95,0.8) {XML};
		\node[papel] at (2.3,-0.95) {adequação da entrada};
		\node[papel] at (5.5,-0.95) {geração da saída};
		% Avaliação: os braços ocupam o lugar da extração e recebem o texto original.
		% Agrupados por quem gera e o que gera; a lista completa fica na Seção 4.6.
		% Os dois grupos que geram DSL ficam em cima, para a saída deles subir até a
		% transpilação sem cruzar a do grupo que gera XML.
		\node[fase] at (-0.1,-1.75) {Avaliação};
		\node[art] (txt2) at (0.4,-3.95) {texto\\bruto};
		\node[prob, text width=3.0cm] (qwen) at (5.5,-2.55) {Modelo pequeno\\gera DSL\\\scriptsize sem e com ajuste fino};
		\node[prob, text width=3.0cm] (frdsl) at (5.5,-3.95) {Modelo grande\\gera DSL\\\scriptsize gramática no \textit{prompt}};
		\node[prob, text width=3.0cm] (frxml) at (5.5,-5.35) {Modelo grande\\gera XML\\\scriptsize geração direta};
		\draw[thick, draw=black!55] (txt2) -- (2.9,-3.95);
		\foreach \b in {qwen, frdsl, frxml} \draw[seta] (2.9,-3.95) |- (\b.west);
		\node[art] at (2.0,-3.65) {texto original};
		% O corpus treina o modelo pequeno (pares texto original e DSL).
		\draw[seta, dotted] (conv.south) -- (8.6,-1.55) -| (qwen.north);
		\node[art, anchor=north west] at (8.7,-0.75) {pares (texto, DSL)\\para o ajuste fino};
		% DSL dos dois grupos converge para a transpilação; XML vai direto à validação.
		\draw[thick, draw=black!55] (qwen.east) -| (7.6,-3.95);
		\draw[thick, draw=black!55] (frdsl.east) -- (7.6,-3.95);
		\fill[black!55] (7.6,-3.95) circle (1.5pt);
		\draw[seta] (7.6,-3.95) -- (11.45,-3.95) -- (trans);
		\node[art] at (10.05,-3.73) {DSL};
		\draw[seta] (frxml.east) -- (14.4,-5.35) -- (val);
		\node[art] at (10.05,-5.13) {XML};
	\end{tikzpicture}
	\Fonte{Elaborado pelo autor}
	\Nota{Borda tracejada indica etapa probabilística; borda sólida, etapa
	determinística. Na construção do corpus, a padronização só adequa o texto de
	entrada, e a extração gera a saída como JSON; daí em diante há apenas
	transformações determinísticas entre representações. Os pares formados pelo
	texto original e pela DSL treinam o modelo pequeno (seta pontilhada). Na
	avaliação, os braços recebem o texto original e ocupam o lugar da extração: os
	que geram DSL passam pela mesma transpilação, e os que geram XML seguem direto
	para a validação.}
\end{figure}
